\documentclass[12pt]{article}

\usepackage[T1]{fontenc}
\usepackage[utf8]{inputenc}
\usepackage[romanian]{babel}
\usepackage{amsmath, amssymb, amsthm}
\usepackage{geometry}
\geometry{a4paper, margin=2.5cm}

\title{Complexitatea Algoritmilor și Teorema Master}
\author{}
\date{}

\begin{document}

\maketitle

\section{Introducere}

Analiza complexității algoritmilor este o parte fundamentală a informaticii.
Scopul ei este să determine cât timp și câtă memorie consumă un algoritm în
funcție de dimensiunea datelor de intrare.

Complexitatea se exprimă în notații asimptotice precum:


\[
O(\cdot),\ \Omega(\cdot),\ \Theta(\cdot),
\]


care descriu comportamentul algoritmului pentru intrări foarte mari.

\section{Complexitatea temporală}

\subsection{Definiție}

Complexitatea temporală a unui algoritm este funcția $T(n)$ care descrie timpul
de execuție în funcție de dimensiunea intrării $n$.

În practică, nu măsurăm timpul real (în secunde), ci numărul de operații
elementare.

\subsection{Notații asimptotice}

\begin{itemize}
    \item $O(f(n))$: limită superioară (cel mult $f(n)$).
    \item $\Omega(f(n))$: limită inferioară (cel puțin $f(n)$).
    \item $\Theta(f(n))$: limită exactă (și $O$, și $\Omega$).
\end{itemize}

\subsection{Cazuri de analiză}

\begin{itemize}
    \item \textbf{Cel mai bun caz}: intrarea este favorabilă.
    \item \textbf{Cazul mediu}: intrarea este aleatoare.
    \item \textbf{Cel mai rău caz}: intrarea este nefavorabilă.
\end{itemize}

\section{Complexitatea spațială}

Complexitatea spațială măsoară memoria suplimentară folosită de algoritm.

\begin{itemize}
    \item Algoritm \textbf{in-place}: folosește $O(1)$ memorie auxiliară.
    \item Algoritm \textbf{out-of-place}: folosește memorie suplimentară $O(n)$.
\end{itemize}

\section{Recursivitate și ecuații de recurență}

Algoritmii recursivi (precum Merge Sort, Quick Sort, Binary Search) sunt analizați
prin \textbf{ecuații de recurență}.

Exemplu:


\[
T(n) = 2T\left(\frac{n}{2}\right) + n.
\]



Aceasta descrie un algoritm care împarte problema în două subprobleme de mărime
$n/2$ și face $n$ operații pentru combinare.

Pentru a rezolva astfel de ecuații folosim \textbf{Teorema Master}.

\section{Teorema Master}

Teorema Master oferă o metodă directă pentru a determina complexitatea recurențelor
de forma:


\[
T(n) = a\,T\left(\frac{n}{b}\right) + f(n),
\]


unde:
\begin{itemize}
    \item $a \ge 1$ este numărul de subprobleme,
    \item $b > 1$ este factorul de împărțire,
    \item $f(n)$ este costul combinării rezultatelor.
\end{itemize}

\subsection{Cazurile Teoremei Master}

Definim:


\[
n^{\log_b a}.
\]



\subsubsection*{Cazul 1: $f(n)$ este mai mic decât $n^{\log_b a}$}

Dacă există $\varepsilon > 0$ astfel încât:


\[
f(n) = O\left(n^{\log_b a - \varepsilon}\right),
\]


atunci:


\[
T(n) = \Theta\left(n^{\log_b a}\right).
\]



\subsubsection*{Cazul 2: $f(n)$ este egal (asimptotic) cu $n^{\log_b a}$}

Dacă:


\[
f(n) = \Theta\left(n^{\log_b a} \log^k n\right),
\]


atunci:


\[
T(n) = \Theta\left(n^{\log_b a} \log^{k+1} n\right).
\]



\subsubsection*{Cazul 3: $f(n)$ este mai mare decât $n^{\log_b a}$}

Dacă există $\varepsilon > 0$ astfel încât:


\[
f(n) = \Omega\left(n^{\log_b a + \varepsilon}\right),
\]


și dacă este îndeplinită condiția de regularitate:


\[
a\,f\left(\frac{n}{b}\right) \le c\,f(n), \quad c < 1,
\]


atunci:


\[
T(n) = \Theta(f(n)).
\]



\section{Exemple de aplicare}

\subsection{Merge Sort}

Răspunde recurenței:


\[
T(n) = 2T\left(\frac{n}{2}\right) + n.
\]



Aici:


\[
a = 2,\quad b = 2,\quad f(n) = n.
\]



Calculăm:


\[
n^{\log_b a} = n^{\log_2 2} = n.
\]



Suntem în \textbf{Cazul 2}, deci:


\[
T(n) = \Theta(n \log n).
\]



\subsection{Binary Search}

Recurența:


\[
T(n) = T\left(\frac{n}{2}\right) + 1.
\]



Aici:


\[
a = 1,\ b = 2,\ f(n) = 1.
\]





\[
n^{\log_b a} = n^{\log_2 1} = n^0 = 1.
\]



Suntem în \textbf{Cazul 2}:


\[
T(n) = \Theta(\log n).
\]



\subsection{Quick Sort (cazul mediu)}

Recurența:


\[
T(n) = 2T\left(\frac{n}{2}\right) + n.
\]



Identică cu Merge Sort:


\[
T(n) = \Theta(n \log n).
\]



\subsection{Quick Sort (cel mai rău caz)}

Recurența:


\[
T(n) = T(n-1) + n.
\]



Aceasta se rezolvă prin sumare:


\[
T(n) = n + (n-1) + \cdots + 1 = \Theta(n^2).
\]



\section{Concluzii}

\begin{itemize}
    \item Analiza complexității este esențială pentru a compara algoritmi.
    \item Notațiile asimptotice descriu comportamentul pentru intrări mari.
    \item Teorema Master este instrumentul principal pentru recurențe divide et impera.
    \item Algoritmi precum Merge Sort și Quick Sort (mediu) au complexitate $O(n\log n)$.
\end{itemize}

\end{document}
